\section{Introdução}\label{sec:introducao}

O presente \emph{template} deve ser seguido para a confecção dos
trabalhos finais de conclusão dos cursos de Ciência da Computação e
Sistemas de Informação da FURB. Ele é fortemente (mas não totalmente)
baseado no modelo da Sociedade Brasileira de Computação de modo a
facilitar posteriores adaptações para publicações das produções dos TCCs
em eventos científicos. O \emph{template} proposto segue as normas da
ABNT para citações e referências bibliográficas, bem como para
referências à figuras, quadros e tabelas.

Para o TCC dos cursos, \textbf{o presente artigo está limitado a 20
páginas}, \ul{incluindo as referências bibliográficas e excluindo os
apêndices e anexos}. Os principais elementos de formatação estão
explicados no anexo no fim deste \emph{template}. \ul{O formato pode ser
diferente caso o artigo seja submetido a um evento científico da área.
Neste caso o estudante pode manter o artigo no formato do evento, mas
deve anexar o comprovante de submissão.} Dúvidas ou problemas devem ser
sanados com a coordenação do TCC do curso.

A introdução deve despertar no leitor o interesse pelo texto,
apresentando os assuntos que serão tratados e o enfoque que será dado ao
tema central. Deve iniciar com uma \textbf{contextualização} do estudo a
ser realizado, explicando claramente sua origem/motivação. A visão geral
do tema deve então ser afunilada até se chegar ao problema a ser
pesquisado. Após o problema ter sido identificado, deve-se delimitar que
aspectos ou elementos serão tratados. Deve-se deixar bem claro o
problema que se quer resolver, com o desenvolvimento do trabalho e
\textbf{justificar} porque o assunto merece ser estudado.

Deve esclarecer a \textbf{formulação do problema} a ser investigado e
deve ser finalizado com os \textbf{objetivos do trabalho} que podem ser
subdivididos \textbf{em geral (obrigatório)} e \textbf{específicos
(opcional)}. O objetivo principal deve ser descrito em uma frase única,
usando o verbo no infinitivo. Os objetivos específicos detalham o
objetivo principal ou definem subprodutos do trabalho. Também se
relacionam a formas de validação ou avaliação do trabalho realizado. Os
objetivos devem ser mensuráveis quanto a se e como foram ou não
atingidos. Os objetivos específicos devem ser enumerados, usando verbos
no infinitivo.

Também deve(m) ser apresentado(s) \textbf{o(s) método(s) de testes} do
trabalho desenvolvido.
